\section{Shadows and end slices in an exponent box}
\label{sec:boxes}

We now replace the monomials of a stable ideal, degree by degree, by
lex segments in the exponent box.  The shadow comparison below will
ensure that the replacement is an ideal.  At the same time we prove
two comparisons for its first and last slices.
The bounded-monomial theorem of Clements and Lindstr\"om gives the shadow
comparison for arbitrary subsets when the bounds are nondecreasing
\cite{CL}.  Here we assume stability of the subsets, allow the bounds in
arbitrary order, and prove the two endpoint comparisons together with
the shadow inequality.

Recall the notation for a box with positive integer bounds
\(b_1,\ldots,b_r\):
\[
 \mathcal B_r(d)=
 \left\{u\in\Z_{\ge0}^{\,r}:
       \sum_{h=1}^r u_h=d,\quad u_h\le b_h\text{ for every }h\right\}.
\]
Here \(d\) may be any integer; in particular the box is empty when
\(d<0\) or \(d>\sum_h b_h\).  We also set
\(\mathcal B_0(0)=\{()\}\) and \(\mathcal B_0(d)=\varnothing\) for
\(d\ne0\).  No ordering of the bounds is required in this section.

Vectors are ordered lexicographically, with the first coordinate
taking precedence, and lex segments start with the largest vectors.
For \(X\subseteq\mathcal B_r(d)\), let
\(\operatorname{Lex}(X)\) be the lex segment of cardinality \(|X|\).
Recall that \(X\) is \emph{stable} if
\[
 u\in X,\quad h<\ell,\quad
 u+e_h-e_\ell\in\mathcal B_r(d)
 \quad\Longrightarrow\quad u+e_h-e_\ell\in X.
\]
Every lex segment is stable, since each allowed transfer increases
the vector in lex order.  Define the bounded shadow and the slices by
\[
 \begin{split}
 \partial X
   &=\{u+e_h\in\mathcal B_r(d+1):u\in X,\ 1\le h\le r\},\\
 X(s)
   &=\{u\in\mathcal B_{r-1}(d-s):(u,s)\in X\}
       \qquad(0\le s\le b_r).
 \end{split}
\]
The bounds in \(\mathcal B_{r-1}\) are the first \(r-1\) bounds.
Empty sets have empty shadows; thus these definitions also cover
negative degrees and empty slices.

\begin{lemma}[Shadows and end slices]
\label{box:inequalities}
Let \(r\ge1\), let \(b_1,\ldots,b_r\) be positive integers, and let
\(d\in\Z\).  If \(X\subseteq\mathcal B_r(d)\) is stable and
\(Y=\operatorname{Lex}(X)\), then \(\partial Y\) is a lex segment and
\begin{equation}
 \label{box:three-inequalities}
 |\partial X|\ge|\partial Y|,
 \qquad |X(0)|\ge|Y(0)|,
 \qquad |X(b_r)|\le|Y(b_r)|.
\end{equation}
\end{lemma}

\begin{proof}
We first prove the assertion about lex shadows, without assuming
stability of \(X\).  Empty shadows need no argument.  For
\(v\in\mathcal B_r(d+1)\), with \(d\ge0\), its lex largest
degree-\(d\) predecessor is
\[
 p(v)=v-e_{\max\{h:v_h>0\}}.
\]
Every predecessor belongs to the bounded box.  Moreover, the map
\(p\) preserves weak lex order.  Indeed, suppose \(v>w\), with first
difference at coordinate \(k\).  Subtraction at the last nonzero
coordinate cannot change an earlier coordinate of \(v\).  If the
operation changes a coordinate of \(w\) before \(k\), then
\(p(v)>p(w)\) already at that coordinate.  Otherwise the strict
inequality at \(k\) can disappear only if \(v_k=w_k+1\) and the tail
of \(v\) after \(k\) is zero.  Equality of total degrees then says
that the tail of \(w\) contains exactly one unit.  Subtracting its
last nonzero coordinate gives \(p(v)=p(w)\).  In all other cases the
first strict inequality remains.  Membership of \(v\) in the shadow
of a lex segment is equivalent to membership of \(p(v)\) in that
segment, proving that its shadow is lex.

We prove the three inequalities simultaneously by induction on
\(r\).  If \(\mathcal B_r(d)\) is empty, they are immediate.  For
\(r=1\), each degree box has at most one vector, so \(X=Y\).
Assume henceforth that \(r>1\).  A prime will indicate lex
replacement or shadow in the first \(r-1\) coordinates.

Each slice \(X(s)\) is stable, and stability between slices gives
\begin{equation}
 \label{box:original-slices}
 \partial'X(s)\subseteq X(s-1)
       \qquad(1\le s\le b_r).
\end{equation}
Replace each slice by \(Z(s)=\operatorname{Lex}'(X(s))\), and let
\(Z\subseteq\mathcal B_r(d)\) be the set with these slices.
By the induction hypothesis and \eqref{box:original-slices},
\[
 |\partial'Z(s)|\le|\partial'X(s)|\le|X(s-1)|=|Z(s-1)|.
\]
The shadow on the left and the slice on the right are lex segments
in the same degree box.  Consequently
\begin{equation}
 \label{box:compressed-slices}
 \partial'Z(s)\subseteq Z(s-1)
       \qquad(1\le s\le b_r).
\end{equation}
Thus a unit in the last coordinate of a member of \(Z\) may be moved
to any earlier coordinate with available capacity, while remaining
in \(Z\).

\smallskip\noindent
\emph{The zero slice.}
If \(Z(0)=\mathcal B_{r-1}(d)\), then
\(|Y(0)|\le|Z(0)|=|X(0)|\).  Otherwise let \(v\) be the largest
vector of \(\mathcal B_{r-1}(d)\) omitted from \(Z(0)\); all members
of \(Z(0)\) are strictly larger than \(v\).

We claim that every \((u,s)\in Z\) is strictly larger than
\((v,0)\).  Equality is excluded by the choice of \(v\).  If instead
\((u,s)<(v,0)\), let \(k<r\) be their first differing coordinate.
Then \(u_k<v_k\) and earlier coordinates agree. Also \(s>0\),
since every member of the lex slice \(Z(0)\) is larger than \(v\).
We will
move the \(s\) units of the last coordinate into the first \(r-1\)
coordinates to obtain a vector \(w\le v\).

Let
\[
 c=\sum_{k<j<r}(b_j-u_j)
\]
be the available capacity strictly after coordinate \(k\).  If
\(s\le c\), distribute all \(s\) units among those positions; the
result still has \(w_k=u_k<v_k\).  If \(s>c\), fill all those
positions and put the remaining \(s-c\) units at coordinate \(k\).
Since \(|u|+s=|v|=d\) and the prefixes before \(k\) agree,
\begin{equation}
 \label{box:tail-filling}
 s-c
   =v_k-u_k-\sum_{k<j<r}(b_j-v_j)
   \le v_k-u_k.
\end{equation}
The remaining units therefore fit at \(k\).  If the inequality in
\eqref{box:tail-filling} is strict, then \(w_k<v_k\).  If equality
holds, every coordinate of the tail of \(v\) is full, and the
constructed vector is exactly \(v\).  In either case \(w\le v\).
All the moves stay within the bounds, so repeated use of
\eqref{box:compressed-slices} gives \(w\in Z(0)\), a contradiction.
This proves the claim.

There are thus at least \(|Z|=|X|\) vectors preceding \((v,0)\).
The lex segment \(Y\) of that cardinality also consists entirely of
vectors preceding \((v,0)\).  Its zero slice is therefore contained
in \(Z(0)\), proving
\begin{equation}
 \label{box:zero-slice}
 |Y(0)|\le|Z(0)|=|X(0)|.
\end{equation}

\smallskip\noindent
\emph{The top slice.}
Put \(B=\sum_{h=1}^r b_h\) and let
\(\rho(u)=(b_1-u_1,\ldots,b_r-u_r)\).  The reflected complement
\[
 X^*=\rho\bigl(\mathcal B_r(d)\setminus X\bigr)
       \subseteq\mathcal B_r(B-d)
\]
is stable.  In fact, if an allowed upward transfer took a member
of \(X^*\) outside \(X^*\), reflection would give an upward transfer
from a member of \(X\) to a vector outside \(X\), contradicting
stability.  Reflection reverses lex order, so
\(Y^*=\rho(\mathcal B_r(d)\setminus Y)\) is the lex segment of
cardinality \(|X^*|\).

The zero-slice inequality, already proved in \(r\) coordinates,
applies to \(X^*,Y^*\).  Their zero-slice cardinalities are
\[
 \begin{split}
 |X^*(0)|&=|\mathcal B_{r-1}(d-b_r)|-|X(b_r)|,\\
 |Y^*(0)|&=|\mathcal B_{r-1}(d-b_r)|-|Y(b_r)|.
 \end{split}
\]
It follows that \(|X(b_r)|\le|Y(b_r)|\).

\smallskip\noindent
\emph{The shadow.}
We have established both end-slice comparisons.  To recover the
shadow bound, we separate its zero slice from the remaining vectors.
For any stable set \(X\),
\begin{equation}
 \label{box:shadow-count}
 |\partial X|=|\partial'X(0)|+|X|-|X(b_r)|.
\end{equation}
Indeed, its shadow vectors with last coordinate zero are precisely
\(\partial'X(0)\).  Every shadow vector \(v\) with positive last
coordinate has \(v-e_r\in X\): if it was obtained as \(v=u+e_h\)
with \(h<r\), transfer one unit of the positive last coordinate of
\(u\) to coordinate \(h\).  This transfer is within bounds because
\(v\) is within bounds.  Stability gives \(v-e_r\in X\).
Thus adding \(e_r\) bijects the members of \(X\) outside its top
slice with shadow vectors having positive last coordinate,
proving~\eqref{box:shadow-count}.

The same identity holds for the stable set \(Y\).  Its zero slice
is a lex segment, and \eqref{box:zero-slice} gives
\(Y(0)\subseteq\operatorname{Lex}'(X(0))\).  The induction
hypothesis and monotonicity of shadows under inclusion imply
\[
 |\partial'X(0)|
 \ge\bigl|\partial'\operatorname{Lex}'(X(0))\bigr|
 \ge|\partial'Y(0)|.
\]
Combining this with the top-slice inequality in
\eqref{box:shadow-count} proves \(|\partial X|\ge|\partial Y|\)
and completes the induction.
\end{proof}

We apply the shadow comparison to the degree pieces of a monomial
ideal.  For \(A=\C[x_1,\ldots,x_r]\) and
\(P_r=(x_1^{b_1+1},\ldots,x_r^{b_r+1})\), a monomial ideal
\(H\supseteq P_r\) is stable in the box sense precisely when each
set
\[
 X_d=\{u\in\mathcal B_r(d):x^u\in H\}
\]
is stable.  Define its degreewise lex replacement by
\begin{equation}
 \label{box:lex-replacement}
 (\Lambda H)_d
   =(P_r)_d+
     \operatorname{span}_{\C}
        \{x^u:u\in\operatorname{Lex}(X_d)\},
 \qquad
 \Lambda H=\bigoplus_{d\ge0}(\Lambda H)_d.
\end{equation}
At this point \(\Lambda H\) is a graded vector subspace; the next
proposition verifies its closure under multiplication.

\begin{proposition}[Lex replacement is an ideal]
\label{box:ideal}
Let \(b_1,\ldots,b_r\) be positive integers, and let
\(H\subseteq\C[x_1,\ldots,x_r]\) be a monomial ideal containing
\(P_r=(x_1^{b_1+1},\ldots,x_r^{b_r+1})\) and stable in the box
sense.  Then the subspace \(\Lambda H\) defined in
\eqref{box:lex-replacement} is a monomial ideal containing \(P_r\).
It is stable in the box sense and has the same Hilbert function as
\(H\).
\end{proposition}

\begin{proof}
Write \(Y_d=\operatorname{Lex}(X_d)\).  Since \(H\) is an ideal,
\(\partial X_d\subseteq X_{d+1}\).  Lemma~\ref{box:inequalities}
therefore gives
\[
 |\partial Y_d|\le|\partial X_d|
                  \le|X_{d+1}|=|Y_{d+1}|.
\]
Both \(\partial Y_d\) and \(Y_{d+1}\) are lex segments in
\(\mathcal B_r(d+1)\), so \(\partial Y_d\subseteq Y_{d+1}\).
Thus multiplying a monomial of \(\Lambda H\) by a variable either
gives a monomial of the next lex segment or leaves the box and
enters \(P_r\).  Products of monomials already in \(P_r\) stay in
\(P_r\).  This proves the ideal assertion.  Each \(Y_d\) is stable
and has cardinality \(|X_d|\), which proves stability and equality
of Hilbert functions.
\end{proof}
